The region map analysis is based on the creation of simple region and then
merges each simple region with relevant ones to build more complex regions.

\subsection{Alias analysis}

\begin{itemize}
    \item modular
    \item over-approximation
    \item flow-insensitive analysis
    \item field/shape analysis
    \item type access analysis (for heterogenous access analysis)
\end{itemize}


\subsection{Shape analysis}
\label{sec-region-shape-analysis}
\subsubsection{Compound merge}

\begin{lstlisting}
struct S {
    int a, b, c[4];
};
void job (int k) {
    struct S s;
    s.a = 1;
    s.c[k] = 1;
}\end{lstlisting}

\begin{table}[h!]
\begin{center}
\begin{tabular}{|m{5cm}|m{10cm}|}
\hline
\centering Compatible compound merge & Meaning
\\
\hline
{\centering{\includegraphics[width=5cm]{"images/example-compound-merge-field-a.png"}}}
&
Regions created by analysing the statement \texttt{s.a = 1}.
\\\hline
{\centering{\includegraphics[width=5cm]{"images/example-compound-merge-field-c.png"}}}
&
Regions created by analysing the statement \texttt{s.c[k] = 1}.
\\
\hline \hline
{\centering{\includegraphics[width=5cm]{"images/example-compound-merge-field-a-c.png"}}}
&
Resulting regions from the merge after analysing both statements.

Since field positions and sizes are compatible among themselves,
the resulting region does not collapses into a garbled region.
\\
\hline
\end{tabular}
\end{center}
\vspace{-15px}
\caption{Examples of region arcs}
\label{table-region-navigation-arcs}
\end{table}
